\documentclass[british]{article}

\usepackage[T1]{fontenc}
\usepackage[latin9]{inputenc}

\usepackage{amstext}
\usepackage{array}
\usepackage{booktabs}
\usepackage{enumitem}
\usepackage{geometry}
\geometry{tmargin=3.5cm,bmargin=3.5cm,lmargin=3cm,rmargin=3cm}
\usepackage{graphicx}
\usepackage{hyperref}
\usepackage{listings}
\lstset{breaklines=true, basicstyle=\ttfamily} 
\usepackage{longtable}
\usepackage{pdflscape}
\usepackage{url}
\usepackage{xcolor}

\usepackage{cleveref}


\newlist{compactenum}{enumerate}{1} % create a custom enumerate-like list
\setlist[compactenum,1]{
	label=\arabic*.,
	leftmargin=*,
	itemsep=0.05mm,
    nolistsep}
\newlist{compactitems}{itemize}{1} % create a custom enumerate-like list
\setlist[compactitems,1]{
	label=\textbullet,
	leftmargin=*,
	itemsep=0.05mm,
    nolistsep}

\newcommand{\PreserveBackslash}[1]{\let\temp=\\#1\let\\=\temp}
\newcolumntype{C}[1]{>{\PreserveBackslash\centering}p{#1}}
\newcolumntype{R}[1]{>{\PreserveBackslash\raggedleft}p{#1}}
\newcolumntype{L}[1]{>{\PreserveBackslash\raggedright}p{#1}}

\title{Fuzzing report for {\texttt @@count_type@@} }
\author{Automatically generated by {\sf $\#$fuzz}}

\begin{document}

\maketitle %

\section{Introduction}
\label{sec:introduction}

{\sf $\#$fuzz} is a fuzzing debugging tool for propositional model counters, developed by Anna L.D. Latour and Mate Soos.
The current version can be found in the (private) repository: \url{https://github.com/meelgroup/count_fuzzer}.

\section{Experimental Setup}
\label{sec:experimental-setup}

\subsection{Fuzzing parameters}
\label{subsec:fuzzing-parameters}

\begin{table}[ht]
    \begin{center}
    \caption{Fuzzing parameters.}
    \label{tab:fuzzing-parameters}
@@table_fuzzing_parameters@@
    \end{center}
\end{table}

\subsection{Instance Generator Information}
\label{subsec:instance-generator-information}

@@list_generator_info@@

\subsection{Counter Information}
\label{subsec:counter-information}

@@list_counter_info@@

\section{Generator results}
\label{sec:generator-results}


\subsection{Satisfiability of generated instances}
\Cref{tab:generator-satisfiability} shows the satisfiability of the instances that were generated by the different instance generators during the fuzzing process.

\begin{table}[ht]
    \begin{center}
    \caption{Satisfiability of instances generated during fuzzing.}
    \label{tab:generator-satisfiability}
@@table_generator_satisfiability@@
    \end{center}
\end{table}

\subsection{Unsatisfiable instances}
\label{subsec:unsatisfiable-instances}

@@list_unsatisfiable_instances@@


\subsection{Instances with unknown satisfiability}
\label{subsec:unknown-instances}

@@list_unknown_instances@@


\section{Counter results}
\label{sec:counter-results}

\begin{table}[ht]
    \begin{center}
    \caption{Aggregate results for counters on all the fuzzing instances. Here `agree' indicates the number of instances for which both the counter and the verifier return a model count and agree on the returned model count. The column `disagree' shows for each counter on how many instances both the counter and the verifier returned a model count, but disagreed on what that model count was. The `counter fails' column indicates on how many instances the counter failed to return a model count (due to time out, memory out, or any other error), while the verifier did successfully return a model count. The `verifier fails' column indicates for how many instances the verifier failed to return a model count, while the counter succeeded. Finally, the `both fail' column indicates for how many instances both the model counter and the verifier failed to return a model count.}
    \label{tab:counter-verifier-status-summary}
@@table_counter_verifier_status_summary@@
    \end{center}
\end{table}

% \begin{table}[ht]
%     \begin{center}
%     \caption{Detailed results for all (counter, instance) pairs for which the fuzzer detected potential problems.}

\newpage
\begin{landscape}
@@table_counter_verifier_status_details@@
\end{landscape}
    % \end{center}
% \end{table}
\end{document}